\documentclass[aspectratio=169]{beamer}
\usepackage[utf8]{inputenc}
\usepackage[T2A]{fontenc}
\usepackage[russian]{babel}
\usetheme[board]{petrovsky}   % [paper] — светлая
\usepackage{pgfplots}\pgfplotsset{compat=1.18}

\title{Производная сложной функции}
\subtitle{Математика, лекция 4}
\author{А. Петровский, ИС-21}
\date{28.09}

\begin{document}
\begin{frame}\titlepage\end{frame}

\section{Правило цепочки}

\begin{frame}{Производная сложной функции}{Лекция 4}
\begin{theorem}{4.2}
Если $g$ дифференцируема в точке $x$, а $f$ — в точке $g(x)$, то
\[ \bigl(f(g(x))\bigr)' = f'(g(x))\cdot g'(x) \]
\end{theorem}
\end{frame}

\begin{frame}{Как меняется наклон}{График}
\begin{tikzpicture}
\begin{axis}[width=12cm,height=5.4cm,axis lines=middle,
  axis line style={pv-muted},tick style={pv-muted},ticklabel style={font=\tiny,pv-muted},
  xmin=-3.3,xmax=3.3,ymin=-1.4,ymax=1.4,samples=200,no markers]
  \addplot[pv-text,line width=1pt,domain=-3.2:3.2]{sin(deg(x^2))};
  \addplot[pv-accent,line width=1pt,domain=-3.2:3.2]{0.25*2*x*cos(deg(x^2))};
\end{axis}
\end{tikzpicture}
\end{frame}

\begin{frame}{Найдите производную}{Задача}
\begin{task}{1}
$y = \sin(3x^2+1)$
\end{task}
\pause
\begin{enumerate}
\item внешняя: $\sin u \;\to\; \cos u$
\item внутренняя: $u = 3x^2+1 \;\to\; u' = 6x$
\end{enumerate}
\pause
\begin{answerblock}
$y' = 6x\cos(3x^2+1)$
\end{answerblock}
\end{frame}
\end{document}
